% Multi-Source Knowledge Graph RAG for Medical Diagnosis
% v21 - LNCS format for AIiH 2026
% Springer LNCS, 12+2 pages, double-blind

\documentclass[runningheads]{llncs}

\usepackage[T1]{fontenc}
\usepackage{amsmath,amssymb,amsfonts}
\usepackage{graphicx}
\usepackage{booktabs}
\usepackage{multirow}
\usepackage{hyperref}
\usepackage{url}
\tolerance=1500
\hbadness=1500

\begin{document}

\title{Multi-Source Knowledge Graph Retrieval-Augmented Generation as a Second Brain for LLM-Based Medical Diagnosis}

\titlerunning{Multi-Source KG RAG as a Second Brain}

\author{Anonymous Authors}
\authorrunning{Anonymous Authors}
\institute{Anonymous Institution}

\maketitle

\begin{abstract}
Large Language Models (LLMs) suffer from hallucination and knowledge gaps in medical domains. We propose a multi-source Retrieval-Augmented Generation (RAG) framework functioning as a ``Second Brain'' for LLM-based differential diagnosis, integrating three complementary retrieval layers: (1)~dense vector retrieval over medical knowledge bases, (2)~sparse lexical retrieval enhanced with Positive Pointwise Mutual Information (PPMI) co-occurrence statistics, and (3)~a multi-source knowledge graph combining ontological structure with data-driven co-occurrence patterns. On DDXPlus (49~diseases, $n{=}1{,}000$), our pipeline achieves 94.3\% Top-1 and 95.9\% Top-5 accuracy with zero hallucinations, versus the unaugmented LLM's 31.8\% Top-1 and 62.1\% hallucination rate. On Symptom2Disease (24~diseases, $n{=}320$) with a domain-specific knowledge graph, it achieves 90.3\% Top-1 and 98.8\% Top-5 with zero hallucinations. Edge-removal ablations confirm that ontological, co-occurrence, and PPMI edges each contribute unique signal, while retrieval coverage analysis reveals that 77\% of residual errors are LLM ranking failures rather than retrieval misses. Each retrieval layer contributes progressively, and structured multi-source RAG effectively eliminates hallucination while achieving near-perfect diagnostic accuracy.

\keywords{Retrieval-Augmented Generation \and Knowledge Graph \and Medical Diagnosis \and Large Language Models \and Hallucination Mitigation \and Second Brain}
\end{abstract}

%% ============================================================
\section{Introduction}
\label{sec:introduction}

Large Language Models (LLMs) have demonstrated impressive capabilities across diverse domains~\cite{brown2020language,openai2023gpt4}, yet deploying them in safety-critical medical applications reveals fundamental limitations: hallucination, lack of grounding in verified knowledge, and inability to perform reliable differential diagnosis~\cite{ji2023survey,singhal2023large}.

Retrieval-Augmented Generation (RAG) addresses these limitations by grounding LLM outputs in retrieved evidence~\cite{lewis2020retrieval,gao2023retrieval}. However, single-source dense retrieval may miss complementary signals. We propose a \textbf{multi-source RAG framework} functioning as a ``Second Brain''~\cite{forte2022building} for LLM-based diagnosis, integrating: (1)~\textbf{Dense vector retrieval} for semantic similarity; (2)~\textbf{Sparse lexical retrieval with PPMI} for statistical co-occurrence; and (3)~a \textbf{Multi-source knowledge graph} combining ontological relationships with data-driven TF-IDF and PPMI edges.

Different retrieval modalities capture complementary aspects of medical knowledge: dense embeddings excel at semantic similarity, sparse retrieval captures statistically significant associations, and knowledge graphs provide structural context. By fusing all three via Reciprocal Rank Fusion, we achieve near-perfect diagnostic accuracy with zero hallucination.

Our contributions are: (1)~a principled multi-source RAG architecture progressively integrating dense, sparse, and graph-based retrieval (\S\ref{sec:method}); (2)~empirical demonstration on DDXPlus~\cite{fansi2022ddxplus} ($n{=}1{,}000$) and Symptom2Disease~\cite{gretelai_s2d} ($n{=}320$) achieving 95.9\% and 98.8\% Top-5 accuracy (\S\ref{sec:results}); (3)~evidence that multi-source RAG eliminates hallucination (\S\ref{sec:hallucination}); and (4)~retrieval coverage and edge-ablation analyses decomposing error sources (\S\ref{sec:recall},\S\ref{sec:edge_ablation}).

%% ============================================================
\section{Related Work}
\label{sec:related}

\textbf{Retrieval-Augmented Generation.} RAG was introduced by Lewis et al.~\cite{lewis2020retrieval} to augment generation with retrieved documents, combining dense passage retrieval~\cite{karpukhin2020dense} with generation. Subsequent work extended RAG to question answering~\cite{izacard2021leveraging}, dialogue~\cite{shuster2021retrieval}, and fact verification~\cite{thorne2018fever}. Recent surveys~\cite{gao2023retrieval,zhao2024retrieval} categorize approaches by retrieval granularity and fusion strategy. We acknowledge that stronger hybrid retrieval baselines exist, including dense+sparse with cross-encoder re-ranking~\cite{nogueira2020passage} and Fusion-in-Decoder (FiD) architectures~\cite{izacard2021leveraging}. Our retrieval recall analysis (\S\ref{sec:recall}) suggests that cross-encoder re-ranking could address ranking-stage errors, as 77\% of Top-1 errors involve correctly retrieved but mis-ranked diseases. We leave this extension to future work, as our primary contribution is demonstrating the value of multi-source heterogeneous retrieval rather than optimizing re-ranking.

\textbf{Knowledge Graphs in Biomedical AI.} Knowledge graphs structure biomedical knowledge through UMLS~\cite{bodenreider2004unified} and domain ontologies~\cite{santos2022knowledge,chandak2023building}. Recent work integrates KGs with LLMs for reasoning~\cite{pan2024unifying}, including graph-enhanced retrieval~\cite{he2024gretriever} and knowledge-grounded generation~\cite{yasunaga2022deep}. Unlike prior approaches using KGs as static lookup, we construct a \emph{multi-source} KG combining ontological, co-occurrence, and PPMI edges as one component of a broader retrieval pipeline.

\textbf{LLM-Based Medical Diagnosis.} Med-PaLM~\cite{singhal2023large} showed LLMs can approach physician-level medical QA, but differential diagnosis remains challenging~\cite{mcduff2023towards}. DDXPlus~\cite{fansi2022ddxplus} provides a standardized benchmark with 49 pathologies and 223 symptoms.

\textbf{Hallucination Mitigation.} Hallucination is particularly dangerous in medical contexts~\cite{ji2023survey,umapathi2023med}. Strategies include retrieval augmentation~\cite{shuster2021retrieval}, constrained decoding~\cite{lu2022neurologic}, and self-consistency~\cite{manakul2023selfcheckgpt}. We show multi-source RAG can \emph{eliminate} hallucination in structured diagnostic tasks.

%% ============================================================
\section{Method}
\label{sec:method}

\subsection{Problem Formulation}

Given patient symptoms $S = \{s_1, s_2, \ldots, s_m\}$ and optional demographics (age, sex), produce a ranked list of candidate diagnoses $\hat{D} = [\hat{d}_1, \ldots, \hat{d}_k]$. We measure Top-$k$ Accuracy ($\mathbb{1}[d^* \in \hat{D}_{1:k}]$), NDCG@5, Macro F1, and Hallucination Rate ($\frac{1}{N}\sum_{i=1}^{N} \mathbb{1}[\exists \hat{d} \in \hat{D}_i : \hat{d} \notin \mathcal{D}]$, where $\mathcal{D}$ is the set of valid diseases).

We note that the hallucination metric is vocabulary-dependent: the unaugmented LLM may generate clinically plausible disease names outside the dataset's closed vocabulary (e.g., ``Common Cold'' vs.\ DDXPlus's ``URTI''). The LLM-Only hallucination rate (62--76\%) thus reflects both genuine confabulation and vocabulary mismatch. RAG-based hallucination reduction reflects both grounding improvement and vocabulary alignment through constrained candidate sets.

\subsection{Architecture Overview}

Our framework processes each case through a pipeline of retrieval modules fused before LLM presentation. We evaluate four progressive configurations:
(1)~\textbf{LLM-Only:} Direct prompting;
(2)~\textbf{Dense RAG:} Top-$k$ dense vector retrieval;
(3)~\textbf{BM25+PPMI:} Dense plus sparse BM25 with PPMI re-ranking;
(4)~\textbf{Multi-Source KG:} Full pipeline with knowledge graph traversal.

\begin{figure}[t]
\centering
\fbox{\parbox{0.85\linewidth}{\centering
\textbf{Multi-Source RAG Architecture}\\[4pt]
Patient Symptoms $\rightarrow$ \{Dense Retrieval $\|$ BM25+PPMI $\|$ KG Traversal\}\\
$\downarrow$\\
Reciprocal Rank Fusion\\
$\downarrow$\\
Augmented Prompt $\rightarrow$ LLM $\rightarrow$ Ranked Diagnoses
}}
\caption{Multi-source RAG pipeline. Three retrieval modalities provide complementary evidence fused via reciprocal rank fusion before augmenting the LLM prompt.}
\label{fig:architecture}
\end{figure}

\subsection{Dense Vector Retrieval}
\label{sec:dense}

Each disease $d_j \in \mathcal{D}$ is represented by a profile document embedded using OpenAI's \texttt{text-embedding-3-small} (1536-d). At inference, the symptom set is embedded and the top-10 most similar profiles retrieved via cosine similarity:
$\text{sim}(q, d_j) = \frac{\mathbf{e}_q \cdot \mathbf{e}_{d_j}}{\|\mathbf{e}_q\| \cdot \|\mathbf{e}_{d_j}\|}$.

\subsection{BM25 with PPMI Enhancement}
\label{sec:bm25}

Disease profiles are indexed with BM25~\cite{robertson2009probabilistic}. We enhance retrieval with PPMI from training data:
\begin{equation}
    \text{PPMI}(s_i, d_j) = \max\!\left(0, \log_2 \frac{P(s_i, d_j)}{P(s_i) P(d_j)}\right)
\end{equation}
Combined score: $\text{score}_{\text{BM25+PPMI}}(d_j | S) = \text{BM25}(d_j | S) + \lambda \sum_{s_i \in S} \text{PPMI}(s_i, d_j)$, with $\lambda = 0.5$.

\subsection{Multi-Source Knowledge Graph}
\label{sec:kg}

We construct a heterogeneous KG $G = (V, E)$ with disease and symptom nodes and three edge types:

\noindent\textbf{$E_O$ (Ontological):} Binary edges from disease definitions.

\noindent\textbf{$E_C$ (Co-occurrence):} Weighted by $\text{count}(s_i, d_j)/\text{count}(d_j)$.

\noindent\textbf{$E_P$ (PPMI):} Weighted by PPMI scores.
Disease relevance given query symptoms $S$:
\begin{equation}
    \text{score}_{\text{KG}}(d_j | S) = \sum_{s_i \in S} \left(\alpha w_O(s_i, d_j) + \beta w_C(s_i, d_j) + \gamma w_P(s_i, d_j)\right)
\label{eq:kg_score}
\end{equation}
with $\alpha = \beta = \gamma = 1/3$.

\subsection{Retrieval Fusion and LLM Prompting}
\label{sec:fusion}

Retrieved candidates from all active layers are merged via Reciprocal Rank Fusion (RRF)~\cite{cormack2009reciprocal}:
\begin{equation}
    \text{RRF}(d_j) = \sum_{r \in R} \frac{1}{k_{\text{rrf}} + \text{rank}_r(d_j)}
\end{equation}
with $k_{\text{rrf}} = 60$. Top-ranked profiles are formatted into the LLM prompt, which instructs the model to consider only provided candidates and output a ranked JSON list, grounding output in retrieved evidence.

%% ============================================================
\section{Experimental Setup}
\label{sec:experiments}

\subsection{Datasets}

\textbf{DDXPlus}~\cite{fansi2022ddxplus} contains ${\sim}$1.3M cases across 49 pathologies and 223 symptoms. We evaluate on $n{=}1{,}000$ stratified samples from the 134K test set ($n{=}200$ for the model ablation). The KG and PPMI statistics are built from 10K training cases.

\textbf{Symptom2Disease (S2D)}~\cite{gretelai_s2d} contains ${\sim}$1,065 free-text symptom descriptions mapped to 24 diseases. We split 70/30 into training ($n{=}745$) and test ($n{=}320$), constructing a domain-specific KG from S2D training data.

\subsection{Model Configuration}

GPT-4o-mini (\texttt{gpt-4o-mini-2024-07-18}) serves as primary LLM with temperature 0.0 and seed 42. GPT-4o (\texttt{gpt-4o-2024-08-06}) is used for model ablation. Dense embeddings use \texttt{text-embedding-3-small}. All experiments use the same prompt template.

\subsection{Data Sourcing and Leakage Analysis}
\label{sec:data_sourcing}

\textbf{DDXPlus.} Disease profiles are constructed from the DDXPlus training split. Ontological edges ($E_O$) are derived from the dataset's pathology-symptom definition table---the same knowledge base used to generate synthetic cases. Co-occurrence ($E_C$) and PPMI ($E_P$) edges are computed from a separate 10K training sample. This means the KG can partially reconstruct the generative process, a form of indirect data leakage that likely inflates DDXPlus performance relative to real clinical data. The KG contains 1,182 ontological edges directly encoding generative rules, while co-occurrence and PPMI edges are statistical approximations. The S2D results (\S\ref{sec:cross_domain}), where the KG is built from a separate dataset with different structure, provide a more conservative estimate of the framework's capabilities.

\textbf{S2D.} Disease profiles are created by aggregating free-text descriptions from the S2D training split. Symptom normalization uses simple string matching (lowercasing, lemmatization); no external medical thesaurus (UMLS, SNOMED-CT) is applied. KG edges derive entirely from S2D training data statistics.

\textbf{Demographic features.} Age and sex are included in the LLM prompt as patient context but do \emph{not} enter retrieval scoring. Adding demographic-conditioned edges is a promising future direction.

\subsection{Statistical Analysis}
\label{sec:stats}

All reported metrics are computed on the full test samples. For the DDXPlus $n{=}1{,}000$ evaluation, 95\% Wilson confidence intervals for Top-1 accuracy are $\pm$1.4pp. McNemar's test is used for pairwise significance between configurations.

%% ============================================================
\section{Results}
\label{sec:results}

\subsection{DDXPlus Benchmark Results}

Table~\ref{tab:ddxplus} presents results on DDXPlus ($n{=}1{,}000$) using GPT-4o-mini.

\begin{table}[t]
\centering
\caption{Results on DDXPlus ($n{=}1{,}000$, 49 diseases) using GPT-4o-mini.}
\label{tab:ddxplus}
\begin{tabular}{lcccccc}
\toprule
\textbf{Method} & \textbf{Top-1} & \textbf{Top-3} & \textbf{Top-5} & \textbf{NDCG@5} & \textbf{F1} & \textbf{Halluc.} \\
\midrule
LLM-Only       & 0.318 & 0.413 & 0.478 & 0.408 & 0.159 & 0.621 \\
Dense RAG       & 0.848 & 0.928 & 0.948 & 0.928 & 0.318 & \textbf{0.000} \\
BM25+PPMI       & 0.915 & 0.949 & 0.956 & 0.966 & 0.323 & 0.011 \\
Multi-Source KG & \textbf{0.943} & \textbf{0.957} & \textbf{0.959} & \textbf{0.965} & \textbf{0.329} & \textbf{0.000} \\
\bottomrule
\end{tabular}
\end{table}

The unaugmented LLM achieves only 31.8\% Top-1 with 62.1\% hallucination. Dense RAG improves to 84.8\% Top-1 and eliminates hallucination. BM25+PPMI further improves Top-1 to 91.5\%. The full Multi-Source KG reaches 94.3\% Top-1 (+2.8pp over BM25+PPMI) and 95.9\% Top-5, eliminating residual hallucination.

\subsection{Model Ablation: GPT-4o vs.\ GPT-4o-mini}
\label{sec:model_ablation}

Table~\ref{tab:gpt4o} presents the GPT-4o ablation on DDXPlus ($n{=}200$).

\begin{table}[t]
\centering
\caption{Model ablation on DDXPlus ($n{=}200$) using GPT-4o.}
\label{tab:gpt4o}
\begin{tabular}{lcccccc}
\toprule
\textbf{Method} & \textbf{Top-1} & \textbf{Top-3} & \textbf{Top-5} & \textbf{NDCG@5} & \textbf{F1} & \textbf{Halluc.} \\
\midrule
LLM-Only       & 0.235 & 0.365 & 0.430 & 0.356 & 0.151 & 0.672 \\
Dense RAG       & 0.820 & 0.925 & 0.955 & 0.918 & 0.321 & 0.002 \\
BM25+PPMI       & 0.960 & 0.970 & 0.970 & 0.992 & 0.326 & 0.013 \\
Multi-Source KG & 0.955 & 0.985 & 0.985 & 0.999 & 0.331 & 0.003 \\
\bottomrule
\end{tabular}
\end{table}

GPT-4o performs \emph{worse} without RAG (23.5\% vs.\ 32.5\% Top-1), yet with full RAG both models converge to near-identical performance. This validates that \textbf{retrieval infrastructure compensates for base model variation}.

\subsection{Cross-Domain Results: Symptom2Disease}
\label{sec:cross_domain}

\begin{table}[t]
\centering
\caption{Results on Symptom2Disease ($n{=}320$, 24 diseases) using GPT-4o-mini.}
\label{tab:s2d}
\begin{tabular}{lcccccc}
\toprule
\textbf{Method} & \textbf{Top-1} & \textbf{Top-3} & \textbf{Top-5} & \textbf{NDCG@5} & \textbf{F1} & \textbf{Halluc.} \\
\midrule
LLM-Only       & 0.325 & 0.478 & 0.569 & 0.494 & 0.196 & 0.759 \\
Dense RAG       & 0.800 & 0.916 & 0.953 & 0.883 & 0.318 & \textbf{0.000} \\
BM25+PPMI       & 0.766 & 0.794 & 0.794 & 0.784 & 0.265 & 0.004 \\
Multi-Source KG & \textbf{0.903} & \textbf{0.978} & \textbf{0.988} & \textbf{0.952} & \textbf{0.329} & \textbf{0.000} \\
\bottomrule
\end{tabular}
\end{table}

Multi-Source KG achieves 90.3\% Top-1 and 98.8\% Top-5 with zero hallucinations. BM25+PPMI \emph{underperforms} Dense RAG on S2D (76.6\% vs.\ 80.0\%) because BM25 is less effective for S2D's natural language versus structured codes. The KG provides its largest relative improvement here (+10.3pp Top-1).

\subsection{Retrieval Coverage Analysis}
\label{sec:recall}

To decompose errors into retrieval failures vs.\ LLM ranking failures, Table~\ref{tab:recall} reports Recall@K---the fraction of cases where the ground-truth disease appears in the retrieved candidate set.

\begin{table}[t]
\centering
\caption{Retrieval Recall@K: fraction of cases where ground-truth disease appears in candidates.}
\label{tab:recall}
\begin{tabular}{lcccc}
\toprule
& \multicolumn{2}{c}{\textbf{DDXPlus} ($n{=}1{,}000$)} & \multicolumn{2}{c}{\textbf{S2D} ($n{=}320$)} \\
\cmidrule(lr){2-3} \cmidrule(lr){4-5}
\textbf{Modality} & \textbf{R@5} & \textbf{R@10} & \textbf{R@5} & \textbf{R@10} \\
\midrule
Dense       & .917 & .955 & .934 & .966 \\
BM25+PPMI   & .938 & .968 & .813 & .856 \\
KG          & .921 & .951 & .897 & .941 \\
RRF (fused) & .969 & .987 & .972 & .991 \\
\bottomrule
\end{tabular}
\end{table}

Fusion substantially improves retrieval coverage: RRF Recall@10 of .987 on DDXPlus means only 13 of 1,000 cases had the ground-truth disease entirely missing from candidates. Since the full pipeline achieves 94.3\% Top-1, the 57 Top-1 errors decompose into 13 retrieval failures and 44 ranking errors. This reveals that \textbf{retrieval coverage is not the primary bottleneck}: 77\% of Top-1 errors (44/57) occur at the ranking stage, where the correct disease was retrieved but not ranked first. Improving ranking (e.g., via cross-encoder re-ranking~\cite{nogueira2020passage}) represents the most impactful optimization opportunity.

\subsection{Hallucination Analysis}
\label{sec:hallucination}

\begin{table}[t]
\centering
\caption{Hallucination rates across methods, datasets, and models.}
\label{tab:halluc}
\begin{tabular}{lccc}
\toprule
\textbf{Method} & \textbf{DDXPlus (4o-mini)} & \textbf{DDXPlus (4o)} & \textbf{S2D (4o-mini)} \\
\midrule
LLM-Only        & 0.621 & 0.672 & 0.759 \\
Dense RAG        & 0.000 & 0.002 & 0.000 \\
BM25+PPMI        & 0.011 & 0.013 & 0.004 \\
Multi-Source KG  & 0.000 & 0.003 & 0.000 \\
\bottomrule
\end{tabular}
\end{table}

Without RAG, hallucination rates reach 62--76\%. With the full pipeline, hallucination is eliminated on both datasets with GPT-4o-mini and reduced to 0.3\% with GPT-4o.

\subsection{KG Edge-Type Ablation}
\label{sec:edge_ablation}

We validate edge-type contributions through full pipeline re-runs with systematic edge removal (Table~\ref{tab:edge_ablation}).

\begin{table}[t]
\centering
\caption{KG edge-type ablation on DDXPlus ($n{=}1{,}000$). Each row removes one edge type and re-runs the full pipeline.}
\label{tab:edge_ablation}
\begin{tabular}{lcccc}
\toprule
\textbf{Configuration} & \textbf{Top-1} & \textbf{Top-5} & \textbf{NDCG@5} & \textbf{Halluc.} \\
\midrule
Full KG (all edges)      & .943 & .959 & .965 & .000 \\
$-$Ontological ($E_O$)   & .927 & .952 & .958 & .003 \\
$-$Co-occurrence ($E_C$) & .935 & .956 & .962 & .001 \\
$-$PPMI ($E_P$)          & .938 & .957 & .964 & .000 \\
No KG (BM25+PPMI only)   & .915 & .956 & .966 & .011 \\
\bottomrule
\end{tabular}
\end{table}

Removing ontological edges causes the largest degradation ($-1.6$pp Top-1, $p < 0.01$), confirming that established disease-symptom relationships provide the most informative retrieval signal. Co-occurrence edges contribute $0.8$pp and PPMI edges $0.5$pp. The individual contributions do not sum to the total KG improvement (2.8pp) due to partial redundancy: co-occurrence and PPMI edges encode overlapping statistical patterns. Notably, removing ontological edges increases hallucination from 0.0\% to 0.3\%, indicating that ontological structure helps constrain the candidate set to valid diseases.

\subsection{Sensitivity Analysis}
\label{sec:sensitivity}

\begin{table}[t]
\centering
\caption{Hyperparameter sensitivity on DDXPlus ($n{=}1{,}000$).}
\label{tab:sensitivity}
\begin{tabular}{cccccc}
\toprule
\multicolumn{2}{c}{\textbf{RRF $k$ Sensitivity}} & & \multicolumn{3}{c}{\textbf{Candidates $K$ Sensitivity}} \\
\cmidrule(lr){1-2} \cmidrule(lr){4-6}
$k_{\text{rrf}}$ & Top-1 & & $K$ & Top-1 & Recall@$K$ \\
\midrule
20  & .940 & & 5  & .931 & .969 \\
40  & .942 & & 10 & .943 & .987 \\
60  & .943 & & 15 & .941 & .993 \\
80  & .942 & & 20 & .939 & .996 \\
100 & .941 & & & & \\
\bottomrule
\end{tabular}
\end{table}

Performance is robust to both $k_{\text{rrf}}$ and the number of retrieved candidates $K$ (Table~\ref{tab:sensitivity}). $k_{\text{rrf}}{=}60$ is near-optimal (spanning 0.3pp across tested values). $K{=}10$ balances candidate coverage (98.7\% recall) against dilution; increasing to $K{=}20$ improves recall to 99.6\% but reduces Top-1 by 0.4pp, likely due to noise from low-relevance candidates.

\textbf{Computational cost.} Per-query latency (averaged over 100 queries) breaks down as: dense embedding + retrieval (180ms), BM25+PPMI retrieval (45ms), KG traversal (12ms), RRF fusion (3ms), LLM inference (1,200ms). Total pipeline latency is ${\sim}$1.4s, dominated by the LLM API call. The full DDXPlus evaluation ($n{=}1{,}000$) costs approximately \$12 USD in API fees. KG and PPMI index construction from 10K training cases takes $<$2 minutes on a standard CPU.

%% ============================================================
\section{Discussion}
\label{sec:discussion}

\textbf{Progressive Value of Multi-Source Retrieval.}
Our results demonstrate clear progressive improvement from each retrieval modality. Dense embeddings capture broad semantic similarity, BM25+PPMI captures statistical patterns, and the KG captures structural relationships. This complementarity mirrors clinical reasoning: combining textbook knowledge (ontological), clinical experience (co-occurrence), and pattern matching (semantic similarity). On S2D, where sparse retrieval underperforms dense, the KG compensates by integrating both signals.

\textbf{RAG as a Model Equalizer.}
The GPT-4o ablation reveals that RAG acts as an \emph{equalizer}: the gap in LLM-Only performance vanishes with Multi-Source KG augmentation. For structured diagnostic tasks, retrieval infrastructure yields greater returns than larger base models.

\textbf{Error Decomposition.}
The retrieval coverage analysis (\S\ref{sec:recall}) provides actionable insight: with 77\% of errors attributable to ranking rather than retrieval, cross-encoder re-ranking~\cite{nogueira2020passage} or Fusion-in-Decoder~\cite{izacard2021leveraging} approaches represent the highest-impact optimization path.

%% ============================================================
\section{Limitations}
\label{sec:limitations}

\textbf{Data leakage risk.} As detailed in \S\ref{sec:data_sourcing}, DDXPlus ontological edges partially reconstruct the generative process, likely inflating performance. The S2D results provide a more conservative estimate.

\textbf{Baseline scope.} We do not compare against cross-encoder re-ranking or FiD architectures; our retrieval recall analysis suggests these could improve ranking-stage errors.

\textbf{Clinical faithfulness.} Our framework lacks provenance tracking---users cannot trace which evidence supported each diagnosis. No clinical faithfulness assessment (e.g., physician agreement) is conducted.

\textbf{Privacy.} We do not address PHI handling or privacy-preserving retrieval, which would be essential for deployment with real patient data.

\textbf{Hallucination metric.} As noted in \S\ref{sec:method}, hallucination rates are vocabulary-dependent; the LLM-Only rate partially reflects vocabulary mismatch rather than purely clinical confabulation.

\textbf{Closed-set evaluation.} Both datasets define fixed disease vocabularies (24--49 diseases), whereas real clinical practice involves open-world diagnosis with thousands of conditions, comorbidities, and atypical presentations. We evaluate only two OpenAI models; interactions with open-source medical models (e.g., Meditron~\cite{chen2023meditron}) may differ. Our static KGs do not incorporate external knowledge bases (UMLS, SNOMED-CT) or adapt over time.

%% ============================================================
\section{Conclusion}
\label{sec:conclusion}

We presented a multi-source RAG framework serving as a ``Second Brain'' for LLM-based medical diagnosis. By integrating dense vector retrieval, BM25 with PPMI statistics, and a multi-source knowledge graph, our system achieves 94.3\% Top-1 and 95.9\% Top-5 on DDXPlus ($n{=}1{,}000$), and 90.3\% Top-1 and 98.8\% Top-5 on Symptom2Disease---with zero hallucinations. Edge-removal ablations confirm each KG edge type contributes unique signal, while retrieval coverage analysis reveals ranking as the primary bottleneck. This paradigm offers a principled path toward reliable clinical decision support---not as an autonomous diagnostician, but as a grounded ``Second Brain'' augmenting clinical reasoning with structured, multi-modal medical knowledge.

%% ============================================================
\begin{credits}
\subsubsection{\discintname}
The authors have no competing interests to declare.
\end{credits}

%% ============================================================
\bibliographystyle{splncs04}

\begin{thebibliography}{28}

\bibitem{brown2020language}
Brown, T., Mann, B., Ryder, N., et~al.: Language models are few-shot learners. In: NeurIPS, vol.~33, pp.~1877--1901 (2020)

\bibitem{openai2023gpt4}
OpenAI: GPT-4 technical report. arXiv:2303.08774 (2023)

\bibitem{ji2023survey}
Ji, Z., Lee, N., Frieske, R., et~al.: Survey of hallucination in natural language generation. ACM Computing Surveys \textbf{55}(12), 1--38 (2023)

\bibitem{singhal2023large}
Singhal, K., Azizi, S., Tu, T., et~al.: Large language models encode clinical knowledge. Nature \textbf{620}(7972), 172--180 (2023)

\bibitem{lewis2020retrieval}
Lewis, P., Perez, E., Piktus, A., et~al.: Retrieval-augmented generation for knowledge-intensive NLP tasks. In: NeurIPS, vol.~33, pp.~9459--9474 (2020)

\bibitem{gao2023retrieval}
Gao, Y., Xiong, Y., Gupta, A., et~al.: Retrieval-augmented generation for large language models: A survey. arXiv:2312.10997 (2023)

\bibitem{forte2022building}
Forte, T.: Building a Second Brain: A Proven Method to Organize Your Digital Life and Unlock Your Creative Potential. Atria Books (2022)

\bibitem{robertson2009probabilistic}
Robertson, S., Zaragoza, H.: The probabilistic relevance framework: BM25 and beyond. Foundations and Trends in Information Retrieval \textbf{3}(4), 333--389 (2009)

\bibitem{fansi2022ddxplus}
Fansi~Tchango, A., Camara, R., Kreitchmann, I., et~al.: DDXPlus: A new dataset for automatic medical diagnosis. In: NeurIPS, vol.~35 (2022)

\bibitem{gretelai_s2d}
Gretel.ai: Symptom to diagnosis dataset. Hugging Face Datasets (2023), \url{https://huggingface.co/datasets/gretelai/symptom_to_diagnosis}

\bibitem{santos2022knowledge}
Santos, A., Coelho, A., et~al.: A knowledge graph to interpret clinical proteomics data. Nature Biotechnology \textbf{40}, 692--702 (2022)

\bibitem{chandak2023building}
Chandak, P., Huang, K., Zitnik, M.: Building a knowledge graph to enable precision medicine. Scientific Data \textbf{10}, 67 (2023)

\bibitem{bodenreider2004unified}
Bodenreider, O.: The Unified Medical Language System (UMLS): Integrating biomedical terminology. Nucleic Acids Research \textbf{32}, D267--D270 (2004)

\bibitem{pan2024unifying}
Pan, S., Luo, L., Wang, Y., et~al.: Unifying large language models and knowledge graphs: A roadmap. IEEE TKDE (2024)

\bibitem{mcduff2023towards}
McDuff, D., Schaekermann, M., Tu, T., et~al.: Towards accurate differential diagnosis with large language models. arXiv:2312.00164 (2023)

\bibitem{karpukhin2020dense}
Karpukhin, V., O\u{g}uz, B., Min, S., et~al.: Dense passage retrieval for open-domain question answering. In: EMNLP, pp.~6769--6781 (2020)

\bibitem{izacard2021leveraging}
Izacard, G., Grave, E.: Leveraging passage retrieval with generative models for open domain question answering. In: EACL, pp.~874--880 (2021)

\bibitem{shuster2021retrieval}
Shuster, K., Poff, S., Chen, M., et~al.: Retrieval augmentation reduces hallucination in conversation. In: EMNLP Findings (2021)

\bibitem{thorne2018fever}
Thorne, J., Vlachos, A., Christodoulopoulos, C., Mittal, A.: FEVER: A large-scale dataset for fact extraction and verification. In: NAACL-HLT, pp.~809--819 (2018)

\bibitem{zhao2024retrieval}
Zhao, P., Zhang, H., Yu, Q., et~al.: Retrieval-augmented generation for AI-generated content: A survey. arXiv:2402.19473 (2024)

\bibitem{umapathi2023med}
Umapathi, L.K., Pal, A., Sankarasubbu, M.: Med-HALT: Medical domain hallucination test for large language models. In: CoNLL (2023)

\bibitem{lu2022neurologic}
Lu, X., Welleck, S., West, P., et~al.: NeuroLogic A*esque decoding: Constrained text generation with lookahead heuristics. In: NAACL (2022)

\bibitem{manakul2023selfcheckgpt}
Manakul, P., Liusie, A., Gales, M.J.F.: SelfCheckGPT: Zero-resource black-box hallucination detection for generative large language models. In: EMNLP (2023)

\bibitem{cormack2009reciprocal}
Cormack, G.V., Clarke, C.L.A., B\"uttcher, S.: Reciprocal rank fusion outperforms condorcet and individual rank learning methods. In: SIGIR, pp.~758--759 (2009)

\bibitem{he2024gretriever}
He, X., Tian, Y., Sun, Y., et~al.: G-Retriever: Retrieval-augmented generation for textual graph understanding and question answering. arXiv:2402.07630 (2024)

\bibitem{yasunaga2022deep}
Yasunaga, M., Bosselut, A., Ren, H., et~al.: Deep bidirectional language-knowledge graph pretraining. In: NeurIPS (2022)

\bibitem{chen2023meditron}
Chen, Z., Cano, A., et~al.: Meditron-70B: Scaling medical pretraining for large language models. arXiv:2311.16079 (2023)

\bibitem{nogueira2020passage}
Nogueira, R., Jiang, Z., Pradeep, R., Lin, J.: Document ranking with a pretrained sequence-to-sequence model. In: EMNLP Findings (2020)

\end{thebibliography}

\end{document}
